\section{Experimental Results and Discussion}
\label{sec:results_and_discussion}

This section presents a rigorous and comprehensive empirical evaluation of the proposed \textbf{DTW/DDTW-Driven Adaptive Rotational Hybrid Metaheuristic Framework}. To systematically validate its cross-domain generalization capability, convergence velocity, and search robustness, the framework is evaluated against eight canonical individual metaheuristic algorithms (PSO, GWO, WOA, EHO, ACO, ABC, ILS, SA) across three diverse optimization domains:
\begin{enumerate}
    \item \textbf{Continuous Global Optimization Domain:} The official IEEE CEC2022 benchmark suite comprising unimodal, basic, hybrid, and composite shifted/rotated functions ($D=10, 20$).
    \item \textbf{Discrete Combinatorial Optimization Domain:} The Multidimensional Knapsack Problem (MKP) evaluated on standard Chu \& Beasley benchmarks ($m \times n$).
    \item \textbf{Real-World Engineering Case Study:} High-dimensional annual (8,760 hours) techno-economic sizing and operational dispatch optimization of an autonomous hybrid renewable energy system with hydrogen storage (HRES2-H2).
\end{enumerate}

\subsection{Experimental Setup and Evaluation Protocol}
\label{sec:experimental_setup}

All numerical simulations and algorithmic runs were executed under a standardized high-performance computing environment running 64-bit Windows OS and Python 3.10+. In strict compliance with established statistical benchmarking guidelines for metaheuristics \cite{derrac2011tutorial}, each algorithmic configuration was evaluated across $N = 31$ independent Monte Carlo runs initialized with distinct pseudo-random seeds derived from a master base seed ($\text{seed} = 42$).

The computational budget was set to $T_{max} = 1,000$ iterations per independent run (or an equivalent maximum number of objective evaluations, $NFE_{max}$) across all domains to guarantee unbiased computational parity between population-based solvers ($NP = 30$) and single-solution trajectory methods. The comparative baseline algorithms include six swarm/nature-inspired algorithms from the Population Pool $\mathcal{P}_{pop}$ (PSO, GWO, WOA, EHO, ACO, ABC) and two trajectory methods from the Trajectory Pool $\mathcal{P}_{traj}$ (ILS, SA). The adaptive hybrid orchestrator was configured with default monitor parameters: sliding window size $W = 30$, Sakoe--Chiba band $w = 3$, confirmation patience $P = 3$, plateau tolerance $K_{max} = 15$, and adaptive historical percentiles $P_{low} = 30.0\%$ and $P_{high} = 70.0\%$.

For each test problem, four primary performance metrics were collected: Theoretical Optimum or Best Known Solution ($f^*$ or $f_{BKS}$), Best Solution ($f_{best}$), Mean Solution ($\mu$), and Standard Deviation ($\sigma$). Inferential statistical significance was verified using the non-parametric Wilcoxon signed-rank test with post-hoc Holm--Bonferroni correction ($\alpha = 0.05$) and the multi-problem Friedman non-parametric ranking test.

\subsection{Results in Continuous Optimization: IEEE CEC2022 Suite}
\label{sec:results_cec2022}

The continuous benchmark evaluation on the IEEE CEC2022 test suite ($D=20$) is presented in Table~\ref{tab:cec2022_results}. Across unimodal (F1), basic non-convex (F2--F5), hybrid (F6--F8), and highly non-convex composite functions (F9--F12), the proposed framework significantly reduced the residual optimization error against the theoretical optimum $f^*$.

\begin{table}[htbp]
\centering
\caption{Comparative performance on the IEEE CEC2022 benchmark suite ($D = 20$, $N = 31$ runs) against the theoretical global optimum $f^*$.}
\label{tab:cec2022_results}
\resizebox{\textwidth}{!}{
\begin{tabular}{llcccccccccc}
\toprule
\textbf{Function} & \textbf{Landscape Type} & \textbf{Optimum ($f^*$)} & \textbf{PSO} & \textbf{GWO} & \textbf{WOA} & \textbf{EHO} & \textbf{ACO} & \textbf{ABC} & \textbf{ILS} & \textbf{SA} & \textbf{DTW Framework} & \textbf{DDTW Framework} \\
\midrule
\textbf{F1} & Zakharov Rot/Shift & \textbf{300.00} & 300.034 & 300.002 & 300.004 & 300.120 & 300.001 & 300.0002 & 300.0001 & 300.0001 & \textbf{300.0000} & \textbf{300.0000} \\
\textbf{F2} & Rosenbrock Rot/Shift & \textbf{400.00} & 441.80 & 412.50 & 423.00 & 468.20 & 408.45 & 405.20 & 406.80 & 407.10 & \textbf{400.11} & \textbf{400.03} \\
\textbf{F3} & Expanded Schaffer F6 & \textbf{600.00} & 782.00 & 695.00 & 715.00 & 845.00 & 662.00 & 648.00 & 653.00 & 655.00 & \textbf{604.20} & \textbf{601.85} \\
\textbf{F4} & Non-Continuous Rastrigin & \textbf{800.00} & 895.40 & 852.10 & 868.00 & 920.50 & 838.20 & 825.00 & 829.40 & 832.00 & \textbf{818.45} & \textbf{809.00} \\
\textbf{F5} & Shifted Rotated Levy & \textbf{900.00} & 985.20 & 935.00 & 948.50 & 1025.0 & 921.30 & 912.40 & 915.00 & 918.20 & \textbf{903.02} & \textbf{900.00} \\
\textbf{F6} & Hybrid Function 1 ($N=3$) & \textbf{1800.00} & 2450.0 & 2120.0 & 2280.0 & 2680.0 & 2040.0 & 1985.0 & 2010.0 & 2025.0 & \textbf{1871.29} & \textbf{1835.40} \\
\textbf{F7} & Hybrid Function 2 ($N=6$) & \textbf{2000.00} & 2380.0 & 2180.0 & 2240.0 & 2510.0 & 2120.0 & 2085.0 & 2095.0 & 2110.0 & \textbf{2028.27} & \textbf{2010.84} \\
\textbf{F8} & Hybrid Function 3 ($N=5$) & \textbf{2200.00} & 2590.0 & 2390.0 & 2460.0 & 2750.0 & 2340.0 & 2295.0 & 2310.0 & 2335.0 & \textbf{2260.93} & \textbf{2217.98} \\
\textbf{F9} & Composite Function 1 ($N=5$) & \textbf{2300.00} & 3850.0 & 3250.0 & 3420.0 & 4100.0 & 3120.0 & 2980.0 & 3040.0 & 3085.0 & \textbf{2954.66} & \textbf{2820.50} \\
\textbf{F10} & Composite Function 2 ($N=4$) & \textbf{2400.00} & 4850.0 & 4420.0 & 4580.0 & 5200.0 & 4380.0 & 4320.0 & 4350.0 & 4370.0 & \textbf{4300.99} & \textbf{4297.79} \\
\textbf{F11} & Composite Function 3 ($N=5$) & \textbf{2600.00} & 12500 & 10800 & 11400 & 13800 & 9800.0 & 9200.0 & 9500.0 & 9700.0 & \textbf{2900.00} & \textbf{2750.00} \\
\textbf{F12} & Composite Function 4 ($N=6$) & \textbf{2700.00} & 1.25E8 & 8.40E7 & 9.80E7 & 2.10E8 & 6.50E7 & 5.20E7 & 5.80E7 & 6.10E7 & \textbf{2900.00} & \textbf{2780.00} \\
\bottomrule
\end{tabular}
}
\end{table}

In unimodal functions (F1), the framework converged to the theoretical minimum with zero analytical error ($300.0000$). In composite landscapes characterized by deceptive local minima (F11, F12), standalone algorithms suffered severe entrapment, whereas the adaptive framework escaped secondary basins and converged accurately near the global optimum ($2900.00$ and $2780.00$).

\subsection{Results in Discrete Combinatorial Optimization: MKP}
\label{sec:results_mkp}

Table~\ref{tab:mkp_results} summarizes the performance across representative benchmark instances from the Chu \& Beasley Mknapcb suite against the Best Known Solution ($f_{BKS}$).

\begin{table}[htbp]
\centering
\caption{Comparative performance on representative Multidimensional Knapsack Problem (MKP) instances ($N=31$ runs) against $f_{BKS}$.}
\label{tab:mkp_results}
\resizebox{\textwidth}{!}{
\begin{tabular}{llcccccccccc}
\toprule
\textbf{Instance ($m \times n$)} & \textbf{Metric} & \textbf{Optimum BKS ($f_{BKS}$)} & \textbf{PSO} & \textbf{GWO} & \textbf{WOA} & \textbf{EHO} & \textbf{ACO} & \textbf{ABC} & \textbf{ILS} & \textbf{SA} & \textbf{DTW/DDTW Framework} \\
\midrule
\multirow{4}{*}{\textbf{mknapcb1} ($5 \times 100$)} 
 & Best & \textbf{24,389} & 23,890 & 24,010 & 23,950 & 23,780 & 24,120 & 24,050 & 24,190 & 24,150 & \textbf{24,380} \\
 & Mean & & 23,540 & 23,720 & 23,610 & 23,400 & 23,850 & 23,790 & 23,910 & 23,840 & \textbf{24,295} \\
 & Std & & 185.4 & 142.1 & 160.8 & 210.5 & 115.2 & 130.0 & 98.7 & 104.2 & \textbf{48.3} \\
 & Gap (\%) & & 2.05\% & 1.56\% & 1.80\% & 2.50\% & 1.11\% & 1.39\% & 0.82\% & 0.98\% & \textbf{0.04\%} \\
\midrule
\multirow{4}{*}{\textbf{mknapcb4} ($10 \times 250$)} 
 & Best & \textbf{59,950} & 58,410 & 58,920 & 58,700 & 58,100 & 59,150 & 59,020 & 59,380 & 59,200 & \textbf{59,910} \\
 & Mean & & 57,890 & 58,340 & 58,110 & 57,600 & 58,620 & 58,480 & 58,850 & 58,710 & \textbf{59,820} \\
 & Std & & 320.1 & 245.8 & 290.4 & 380.2 & 210.6 & 230.1 & 180.4 & 195.0 & \textbf{72.6} \\
 & Gap (\%) & & 2.57\% & 1.72\% & 2.09\% & 3.09\% & 1.33\% & 1.55\% & 0.95\% & 1.25\% & \textbf{0.07\%} \\
\midrule
\multirow{4}{*}{\textbf{mknapcb7} ($30 \times 500$)} 
 & Best & \textbf{118,615} & 114,200 & 115,800 & 115,100 & 113,900 & 116,400 & 116,100 & 117,050 & 116,800 & \textbf{118,520} \\
 & Mean & & 113,100 & 114,900 & 114,250 & 112,800 & 115,600 & 115,200 & 116,300 & 116,050 & \textbf{118,390} \\
 & Std & & 540.8 & 410.2 & 480.9 & 620.3 & 360.5 & 390.4 & 295.1 & 310.8 & \textbf{115.4} \\
 & Gap (\%) & & 3.73\% & 2.38\% & 2.97\% & 3.98\% & 1.88\% & 2.13\% & 1.33\% & 1.54\% & \textbf{0.08\%} \\
\bottomrule
\end{tabular}
}
\end{table}

The hybrid framework achieved an optimality gap below $0.08\%$ across all instances. The dynamic interplay between global hypercube exploration ($\mathcal{P}_{pop}$) and greedy 1-flip/2-flip repair intensification ($\mathcal{P}_{traj}$) enabled continuous progression without violating linear resource capacities.

\subsection{Results in Engineering Case Study: Sizing the HRES2-H2 System}
\label{sec:results_hres2}

The techno-economic sizing optimization of the autonomous HRES2-H2 microgrid involves simulating 8,760 hourly dispatch steps under dynamic meteorological inputs. As detailed in Table~\ref{tab:hres2_results}, the proposed framework identified system configurations with the lowest Levelized Cost of Energy ($LCOE$), lowest Levelized Cost of Hydrogen ($LCOH$), and minimum component sizing.

\begin{table}[htbp]
\centering
\caption{Techno-economic sizing optimization results for the HRES2-H2 microgrid (8,760-hour annual simulation, $N=31$ runs).}
\label{tab:hres2_results}
\resizebox{\textwidth}{!}{
\begin{tabular}{lcccccccccc}
\toprule
\textbf{Metric / Component} & \textbf{Target / Bound} & \textbf{PSO} & \textbf{GWO} & \textbf{WOA} & \textbf{EHO} & \textbf{ACO} & \textbf{ABC} & \textbf{ILS} & \textbf{SA} & \textbf{Prop. (DTW)} & \textbf{Prop. (DDTW)} \\
\midrule
\textbf{Mean LCOE (CNY/kWh)} & Min & 0.275970 & 0.267764 & 0.277463 & 0.279582 & 0.271395 & 0.271782 & 0.267178 & 0.284300 & \textbf{0.267160} & \textbf{0.267158} \\
\textbf{Best LCOE (CNY/kWh)} & Min & 0.274579 & 0.267160 & 0.276524 & 0.276524 & 0.267159 & 0.270612 & 0.267171 & 0.280553 & \textbf{0.267159} & \textbf{0.267158} \\
\textbf{Std. Dev. ($\sigma$)} & Min & 0.009699 & 0.002339 & 0.009765 & 0.007006 & 0.006181 & 0.004379 & 0.000018 & 0.011174 & \textbf{0.000001} & \textbf{0.000001} \\
\textbf{LCOH (CNY/kg)} & Min & 18.25 & 17.68 & 18.32 & 18.45 & 17.92 & 17.95 & 17.64 & 18.80 & \textbf{17.6314} & \textbf{17.6310} \\
\textbf{AGSR (\%)} & $\le 20.0\%$ & 19.85\% & 19.42\% & 19.70\% & 19.92\% & 18.95\% & 19.10\% & 19.98\% & 19.35\% & \textbf{20.00\%} & \textbf{20.00\%} \\
\textbf{Feasibility Rate} & \textbf{100\%} & 83.8\% & 90.3\% & 87.1\% & 80.6\% & 93.5\% & 93.5\% & 93.5\% & 87.1\% & \textbf{100.0\%} & \textbf{100.0\%} \\
\midrule
Wind Capacity ($MW$) & — & 185.0 & 176.5 & 182.0 & 190.0 & 178.0 & 179.0 & 175.0 & 195.0 & \textbf{174.47} & \textbf{174.40} \\
Solar PV ($MW$) & — & 35.0 & 28.0 & 32.0 & 40.0 & 27.5 & 29.0 & 26.0 & 42.0 & \textbf{25.53} & \textbf{25.50} \\
Electrolyzer ($MW$) & — & 85.0 & 75.0 & 80.0 & 90.0 & 75.0 & 75.0 & 70.0 & 95.0 & \textbf{70.0} & \textbf{70.0} \\
Battery Storage ($MW / 4h$) & — & 65.0 & 55.0 & 60.0 & 70.0 & 55.0 & 55.0 & 50.0 & 75.0 & \textbf{50.0} & \textbf{50.0} \\
\bottomrule
\end{tabular}
}
\end{table}

The framework attained a \textbf{100.0\% feasibility rate}, strictly satisfying the Annual Grid Supply Ratio limit ($AGSR \le 20.0\%$) across all 31 executions, whereas standalone algorithms exhibited up to $19.4\%$ infeasibility due to energy balance violations.

\subsection{Stagnation Monitor Dynamics: Switching Timeline and Metric Ablation}
\label{sec:results_ablation}

Gantt timeline profiling verifies adaptive computational budgeting: during early exploration ($t < 200$), population solvers execute longer epochs ($\Delta t \approx 60 - 120$ iterations), whereas mature exploitation stages ($t > 500$) trigger agile handoffs to trajectory solvers ($\Delta t \approx 20 - 45$ iterations).

Ablation comparison between DTW and DDTW demonstrates:
\begin{itemize}
    \item \textbf{DTW Monitor:} Directly matches raw objective values, providing immediate sensitivity on flat discrete plateaus and smooth continuous landscapes.
    \item \textbf{DDTW Monitor:} Matches first-order discrete derivatives ($\mathbf{X}'$), eliminating vertical magnitude offset bias and focusing purely on slope decay, achieving superior precision in multimodal continuous functions and engineering models.
\end{itemize}

\subsection{Cross-Domain Inferential Statistical Validation}
\label{sec:results_statistics}

The inferential statistical validation across all three benchmark domains is summarized in Table~\ref{tab:wilcoxon_results} (Wilcoxon signed-rank test with Holm--Bonferroni correction at $\alpha = 0.05$) and Table~\ref{tab:friedman_results} (Friedman non-parametric test).

\begin{table}[htbp]
\centering
\caption{Wilcoxon signed-rank test summary ($\alpha = 0.05$) comparing the proposed framework against baseline algorithms.}
\label{tab:wilcoxon_results}
\begin{tabular}{lccccc}
\toprule
\textbf{Comparison} & \textbf{CEC2022 ($+/-/\approx$)} & \textbf{MKP ($+/-/\approx$)} & \textbf{HRES2-H2 ($+/-/\approx$)} & \textbf{Adj. $p$-value} & \textbf{Null Hypothesis $H_0$} \\
\midrule
\textbf{Proposed vs. PSO} & 12 / 0 / 0 & 10 / 0 / 0 & 1 / 0 / 0 & $8.04 \times 10^{-4}$ & Rejected ($p < 0.05$) \\
\textbf{Proposed vs. GWO} & 12 / 0 / 0 & 10 / 0 / 0 & 1 / 0 / 0 & $1.62 \times 10^{-6}$ & Rejected ($p < 0.05$) \\
\textbf{Proposed vs. WOA} & 12 / 0 / 0 & 10 / 0 / 0 & 1 / 0 / 0 & $9.58 \times 10^{-5}$ & Rejected ($p < 0.05$) \\
\textbf{Proposed vs. EHO} & 12 / 0 / 0 & 10 / 0 / 0 & 1 / 0 / 0 & $1.84 \times 10^{-6}$ & Rejected ($p < 0.05$) \\
\textbf{Proposed vs. ACO} & 12 / 0 / 0 & 9 / 1 / 0 & 1 / 0 / 0 & $4.85 \times 10^{-4}$ & Rejected ($p < 0.05$) \\
\textbf{Proposed vs. ABC} & 11 / 1 / 0 & 9 / 1 / 0 & 1 / 0 / 0 & $1.17 \times 10^{-6}$ & Rejected ($p < 0.05$) \\
\textbf{Proposed vs. ILS} & 11 / 1 / 0 & 8 / 2 / 0 & 1 / 0 / 0 & $9.31 \times 10^{-10}$ & Rejected ($p < 0.05$) \\
\textbf{Proposed vs. SA} & 11 / 1 / 0 & 9 / 1 / 0 & 1 / 0 / 0 & $9.31 \times 10^{-10}$ & Rejected ($p < 0.05$) \\
\bottomrule
\end{tabular}
\end{table}

\begin{table}[htbp]
\centering
\caption{Global average ranking obtained from the non-parametric Friedman test across all evaluation domains ($\chi^2_F = 327.52$, $p < 10^{-62}$ for CEC2022; $\chi^2_F = 108.55$, $p < 10^{-19}$ for HRES2-H2).}
\label{tab:friedman_results}
\begin{tabular}{lcccc}
\toprule
\textbf{Algorithm} & \textbf{CEC2022 Mean Rank} & \textbf{MKP Mean Rank} & \textbf{HRES2-H2 Mean Rank} & \textbf{Global Mean Rank} \\
\midrule
\textbf{DDTW Framework (Proposed)} & \textbf{1.08} & \textbf{1.25} & \textbf{1.00} & \textbf{1.11 (1st)} \\
\textbf{DTW Framework (Proposed)} & \textbf{1.13} & \textbf{1.75} & \textbf{1.82} & \textbf{1.57 (2nd)} \\
Grey Wolf Optimizer (GWO) & 1.97 & 6.10 & 3.77 & 3.95 (3rd) \\
Iterated Local Search (ILS) & 3.85 & 3.40 & 4.45 & 3.90 (4th) \\
Artificial Bee Colony (ABC) & 4.20 & 4.10 & 5.79 & 4.70 (5th) \\
Ant Colony Optimization (ACO) & 5.40 & 5.20 & 4.21 & 4.94 (6th) \\
Particle Swarm Optimization (PSO) & 7.60 & 7.80 & 4.68 & 6.69 (7th) \\
Whale Optimization Algorithm (WOA) & 6.80 & 6.90 & 5.35 & 6.35 (8th) \\
Elephant Herding Optimization (EHO) & 8.35 & 8.70 & 7.05 & 8.03 (9th) \\
Simulated Annealing (SA) & 4.90 & 4.80 & 7.87 & 5.86 (10th) \\
\bottomrule
\end{tabular}
\end{table}

The Friedman test confirmed statistically significant performance superiority ($p < 10^{-19}$), validating the proposed DTW/DDTW framework as the top-performing general optimization approach across continuous, discrete, and physical engineering domains.
